\section{Conclusion}
\label{sec:conclusion}

We presented a knowledge-based belief-space planning framework that integrates proactive human querying into the robot's state-estimation and planning loop. The robot maintains a committed symbolic knowledge state, constructs an action-conditioned frontier belief after execution, updates particle weights from observations, and uses entropy-normalized confidence to decide whether the belief can be committed back to the knowledge base. When confidence is insufficient, the system asks a targeted state-level question selected by expected posterior entropy. This formulation treats human feedback as belief refinement rather than action-level control. Across waste sorting and tomato harvesting, the results show that belief-triggered querying can preserve high task success while avoiding the interaction cost of always querying. The user study provides secondary evidence that state-level querying can support awareness and reduce perceived burden relative to action-level interaction baselines.
